\section{Contoh}
\label{sec:contoh-bab12}

Dua contoh berikut memperlihatkan penerapan mode pengalamatan.
Contoh pertama adalah program utuh yang mengakses array dengan tiga
mode berbeda. Contoh kedua membandingkan ketujuh mode untuk mengakses
elemen yang sama.

\subsection{Contoh 1: Tiga Mode untuk Array yang Sama}

Program pada Kode~\ref{lst:mode-array} menjumlahkan lima elemen array
dengan tiga cara: \textit{direct memory}, \textit{direct index}, dan
\textit{base relative}.

\lstinputlisting[
    language=MASM,
    caption={Program \file{mode-array.asm} membandingkan tiga mode
    pengalamatan.},
    label={lst:mode-array}
]{\CodeDir/mode-array.asm}

Bagian penting program:

\begin{itemize}
    \item Cara pertama memakai \textit{direct memory}:
          \lstasm|MOV AX, data1+0| sampai \lstasm|ADD AX, data1+8|.
          Offset ditulis sebagai konstanta.
    \item Cara kedua memakai \textit{direct index}:
          \lstasm|ADD AX, data1[SI]|. \reg{SI} dinaikkan dua setiap
          elemen karena satu word berukuran dua byte.
    \item Cara ketiga memakai \textit{base relative}:
          \lstasm|ADD AX, [BX+SI]|. \reg{BX} menunjuk awal array dan
          \reg{SI} menjadi indeks.
    \item Ketiga hasil dibandingkan dengan \asm{CMP}. Bila sama,
          program mencetak \code{Ketiga hasil identik.}
\end{itemize}

Keluaran program tampak pada Kode~\ref{lst:keluaran-mode}.

\begin{lstlisting}[
    style=sessiongaya,
    caption={Keluaran program \file{mode-array.asm}.},
    label={lst:keluaran-mode}
]
Perbandingan tiga mode pengalamatan
Ketiga hasil identik.
\end{lstlisting}

\subsection{Contoh 2: Tujuh Mode untuk Satu Elemen}

Misalkan array \code{data1 DW 10,20,30,40,50} dan elemen yang dituju
adalah elemen ke-2 (indeks 2, offset 4, nilai 30). Ketujuh mode dapat
dipakai seperti pada Tabel~\ref{tab:tujuh-mode-contoh}.

\begin{table}[H]
    \centering
    \caption{Tujuh mode pengalamatan untuk mengakses
    \code{data1[2]}.}
    \label{tab:tujuh-mode-contoh}
    \small
    \begin{tabular}{@{}p{2.6cm}p{6.2cm}p{4.2cm}@{}}
        \toprule
        \textbf{Mode} & \textbf{Potongan program} &
        \textbf{Alamat efektif} \\
        \midrule
        \textit{Register} & \lstasm|MOV AX, BX| &
        Tidak ada akses memori \\
        \addlinespace
        \textit{Immediate} & \lstasm|MOV AX, 30| &
        Tidak ada akses memori \\
        \addlinespace
        \textit{Direct} & \lstasm|MOV AX, data1+4| &
        offset \code{data1} $+ 4$ \\
        \addlinespace
        \textit{Indirect} & \lstasm|LEA SI, data1+4| lalu
        \lstasm|MOV AX, [SI]| & offset \code{data1} $+ 4$ \\
        \addlinespace
        \textit{Base relative} & \lstasm|LEA BX, data1| lalu
        \lstasm|MOV AX, [BX+4]| & offset \code{data1} $+ 4$ \\
        \addlinespace
        \textit{Direct index} & \lstasm|LEA SI, data1| lalu
        \lstasm|MOV AX, [SI+4]| & offset \code{data1} $+ 4$ \\
        \addlinespace
        \textit{Based indexed} & \lstasm|LEA BX, data1|,
        \lstasm|MOV SI, 4|, lalu \lstasm|MOV AX, [BX+SI]| &
        offset \code{data1} $+ 4$ \\
        \bottomrule
    \end{tabular}
\end{table}

Kelima mode memori menghasilkan nilai yang sama, yaitu 30, karena
alamat efektifnya sama. Yang berbeda adalah cara membentuk alamat
tersebut. Mode \textit{register} dan \textit{immediate} tidak
mengakses memori sama sekali.

\begin{infobox}
\vspace{1.1em}
Pada mode \textit{direct index}, perpindahan ditulis di dalam
kurung siku bersama register indeks, misalnya \lstasm|[SI+4]|.
Pada mode \textit{indirect}, kurung siku hanya berisi register,
misalnya \lstasm|[SI]|. Perbedaan kecil ini menentukan apakah
alamat efektif ditambah perpindahan atau tidak.
\end{infobox}
